% Companion Part VI -- VI/34 editorial bridge. Addition only; no problem replacements.
\paragraph{A grading remembers algebraic scale.}
An \(\mathbb N\)-graded ring is a ring \(S\) decomposed as
\[
 S=\bigoplus_{d\geq 0} S_d,\qquad S_dS_e\subseteq S_{d+e}.
\]
A nonzero element of \(S_d\) is \emph{homogeneous of degree} \(d\).
For the polynomial ring \(S=k[x_0,x_1,x_2]\) with each variable of
degree one, \(x_0x_2\) and \(x_1^2\) are homogeneous of degree two,
whereas \(x_0+x_1^2\) is not homogeneous. The grading separates
equations that transform consistently when all coordinates are
rescaled. This is the algebraic feature that makes projective geometry
possible.

\paragraph{Homogeneous ideals and the projective zero locus.}
An ideal \(I\) is homogeneous when it is generated by homogeneous
elements, equivalently when every homogeneous component of an element
of \(I\) again belongs to \(I\). In that case, \(S/I\) inherits a
grading: \((S/I)_d\) is the image of \(S_d\). Consider
\[
 I=(x_1^2-x_0x_2)\subseteq k[x_0,x_1,x_2].
\]
Both terms of its generator have degree two, so the quotient is
graded. Its projective zero locus is a conic in \(\mathbb P^2_k\).
By contrast, \((x_0+x_1^2)\) is not homogeneous for the standard
grading; its equation depends on the choice of scale and cannot
be used unchanged as a projective homogeneous equation.

\paragraph{The irrelevant ideal is not an arbitrary exclusion.}
For a positively graded ring, set
\[
 S_+=\bigoplus_{d>0}S_d.
\]
In the standard polynomial ring it is \((x_0,\ldots,x_n)\).
The homogeneous prime \(S_+\) is the vertex of the affine cone.
Projective points must have at least one nonzero homogeneous coordinate,
so \(\operatorname{Proj}S\) excludes homogeneous primes containing
\(S_+\). In a standard graded polynomial ring, this ensures that
the chart opens \(D_+(x_i)\) cover all of \(\operatorname{Proj}S\).
Do not replace this condition by ``all nonzero prime ideals'';
projective points may correspond to a homogeneous prime such as
\((0)\), and in a general graded ring the irrelevant-ideal condition
is the correct definition.

\paragraph{Degree-zero localization removes the ambiguity of scale.}
Let \(f\in S_d\) be homogeneous with \(d>0\). In \(S_f\), a
homogeneous fraction \(a/f^m\) has degree \(\deg a-md\).
The subring of degree-zero fractions is
\[
 (S_f)_0=\left\{\frac{a}{f^m}:a\in S_{md},\ m\geq0\right\}.
\]
If \(S=k[x_0,x_1]\) and \(f=x_0\), then
\[
 (S_{x_0})_0=k\!\left[\frac{x_1}{x_0}\right]\cong k[t].
\]
Although \(x_0^{-1}\) exists in the full localization \(S_{x_0}\),
it has degree \(-1\) and therefore does \emph{not} belong to
\((S_{x_0})_0\). This distinction between the full localization
and its degree-zero subring is indispensable in the \(\operatorname{Proj}\)
construction.
